% =====================================================================
%  Chapter 36 — When the Machine Is Wrong
% =====================================================================

\chapter{When the Machine Is Wrong}\idx{Wrong}
\label{ch:wrong}

Everything in this book has been about a machine that does almost nothing.
This chapter is about the other case: what happens when a system like this
one is wrong, and what the reader should do when it is.

\section{Three ways it fails}

A forecasting layer fails in three distinguishable ways, and they have
different signatures and different remedies.

\textbf{It is wrong about the sign.} The model predicts that a fund will
outperform its peers and it underperforms. The remedy is the obvious one:
the tilt is bounded, the losses are small, and the next month produces a
new forecast. This is the failure mode\idx{Failure mode} the architecture was designed for,
and it is why the tilt's contribution is a fraction of a basis point
rather than a catastrophe.

\textbf{It is right about the sign and wrong about the size.} The model
says a fund will outperform and it does, but by less than the forecast.
The cost is the same as above: the tilt was small.

\textbf{It is right about the sign and wrong about the horizon.} The model
says a fund will outperform over three months and it underperforms over
twelve. This is the failure mode the ensemble's two horizons and the
discounted score\idx{Discounted score} are designed to catch, and Chapter 19 showed that the
ensemble's members routinely disagree.

None of these three is dangerous, and that is the point of the
architecture.

\section{The failure that is dangerous: the regime}\idx{Regime}

The dangerous failure is not about any particular forecast. It is about
the relationship between the model and the market changing. This is where
the tail risk\idx{Tail risk} lives: the three failures above cost fractions of a basis
point, and this one is the failure that the architecture cannot price.

Three mechanisms would do it. A change in the composition of the investor
base --- the crowding mechanism of Chapter 3 --- would compress the
returns that the features were picking up. A change in the structure of
the market itself --- the arrival of a new instrument, a change in the
index construction rules, a change in the tax code --- would break the
mapping from the features to the returns. And a change in the
distribution of returns --- a crisis, a regime of much higher volatility
--- would change the scale on which every threshold\idx{Threshold} in the system sits.

The system's defences against this are partial and worth listing
honestly. The discounted ensemble score is a defence of the first kind:
a member whose edge has decayed stops being weighted. The half-life is a
parameter precisely so that the speed of the defence can be argued about
rather than assumed. The drawdown brake is a defence of the third kind,
and the shock of a regime change is precisely the sort of thing that
produces the drawdown the brake responds to.

What the system has no defence against is the second kind. A change in
the market's structure does not announce itself in the model's own
performance until after it has happened.

\begin{takeaway}[The honest limit]
The system can notice that a forecast has stopped working, because the
ensemble's score is computed from realised outcomes and is discounted. It
cannot notice that the feature set has stopped describing the market,
because that is a change in the relationship rather than in the
performance of any particular forecast.
\end{takeaway}

\section{What a reader should watch}

Four observables, all of them available from the record without any
additional machinery.

\textbf{The tilt's contribution.} Chapter 34 computes it over every
starting month. A reader who runs the system can compute the same thing
each year and watch whether the tilt's contribution has turned
persistently negative. If it has, the model has stopped working, and the
honest response is to turn it off --- which is one of the book's
ablations and requires no new code.

\textbf{The ensemble's weight concentration.} If one member's weight
collapses towards one, the model has found a single configuration it
believes in, which is both a good sign and a warning: it means the
evidence is being carried by one member, and the between-model uncertainty
that the architecture relies on has gone.

\textbf{The mean cash weight.} A rise in the cash weight means the system
is finding less edge, and it is the system's own word on the matter.

\textbf{The drawdown brake's firing frequency.} A brake that fires
constantly is either in a regime the system was not designed for or has a
threshold that is wrong for the reader's circumstances.

The ensemble's weights over the record are shown in \figref{fig:model-weights}, and a reader should watch them for concentration.

\section{When the reader should override the system}\idx{Override}

This is the question Chapter 8 is really about, and the book should
answer it directly.

The system should be overridden when the reader's circumstances have
changed --- not when the portfolio has. A change in income, in the
stability of employment, in the horizon, or in the willingness to sit
through a drawdown is a legitimate reason to change the strategic
allocation, and the book's architecture permits exactly that, because the
allocation was never the model's to set.

The system should not be overridden because the portfolio is down. That is
the disposition effect of Chapter 9 and the myopic loss aversion of
Chapter 8, and both of them are worth more to the investor than the
system is.

And the system should not be overridden because a particular month's
forecast looks wrong to the reader. The forecast is a posterior mean with
a large uncertainty attached, the position it produced was small for
exactly that reason, and the reader's opinion of a specific fund is not
more informative than the model's.

\begin{keytakeaways}
\item A forecasting layer that is wrong about the sign, the size or the
horizon is dangerous only in proportion to the size of the positions it
produced, which is why the tilt is bounded and priced.
\item The dangerous failure is a regime change, and the system's defences
against it are partial: a discounted score notices a decaying edge, but
nothing notices a change in the mapping from features to returns.
\item Four observables let a reader know the model has stopped working: the
tilt's contribution, the ensemble's weight concentration, the cash
weight, and the brake's firing frequency.
\item Override the system when your circumstances change; do not override
it because the portfolio is down.
\end{keytakeaways}
